\section{Experiment 1: Adaptation}
\label{sec:safety}

Of the 46,640 returned replies, 9,405 are refusals and 37,235 compliances, so the
panel refuses 20.2\% of what it returned. Of the 200 scenarios, 125 are Rights or
Benign and expect compliance at every age, 50 are Harmful and expect refusal at
every age, and 25 are Age Restricted and are the only ones whose expected answer
moves.

Two questions run through this section and the measures that answer them are kept
apart. Whether behaviour moves with the disclosed age is asked of Refusal Rate.
Whether the behaviour was right is asked of Action Alignment and the two
directional diagnostics. More refusal is not thereby safer: a model that refuses
a rights-seeking request has conditioned on age and failed. Supplementary tables
are in Appendix~\ref{app:safety}.

\subsection{Safety}
\label{sec:safety-outcomes}

% Scenario-specific safety rates, one panel a model.
% Chapter 4, Section 4.2.1.
%
% Drawn by scripts/figuresafe.py. Rates are scenario weighted through
% analysis.rate_by_model, so a scenario counts once whatever it returned and the
% panel totals agree with Table~\ref{tab:safety-outcome}.
%
% Include at \textwidth and do not scale.
\begin{figure}[htbp]
\centering
\includegraphics[width=0.8\textwidth]{safe/safety_rates.pdf}
\caption{\emph{Experiment 1} $\vert$ Safety Rates by Scenario Type.}
\label{fig:safety-rates}
\end{figure}



\begin{table}[htbp]
\centering
\footnotesize
\setlength{\tabcolsep}{4pt}
\renewcommand{\arraystretch}{1.15}
\begin{tabularx}{\textwidth}{
    @{}
    L{3.9cm}
    >{\raggedleft\arraybackslash}X
    >{\raggedleft\arraybackslash}X
    >{\raggedleft\arraybackslash}X
    >{\raggedleft\arraybackslash}X
    >{\raggedleft\arraybackslash}X
    >{\raggedleft\arraybackslash}X
    >{\raggedleft\arraybackslash}X
    @{}
}
\toprule
 & \multicolumn{4}{c}{\textbf{Outcome (\%)}} & \multicolumn{3}{c}{\textbf{Derived (\%)}} \\
\cmidrule(lr){2-5}\cmidrule(l){6-8}
\textbf{Model} & \textbf{Strong Refusal} & \textbf{Weak Refusal} &
\textbf{Minimal Compliance} & \textbf{Total Compliance} & \textbf{Refusal} &
\textbf{Delivery} & \textbf{Mismatch} \\
\midrule
\modeltag{openaiPastel}{GPT-5.6 Luna}          & 19.2 & 0.2 & 2.7 & 77.9 & 19.4 & 78.2 & 2.9 \\
\modeltag{anthropicPastel}{Claude Haiku 4.5}   & 18.8 & 0.1 & 4.7 & 76.4 & 19.0 & 76.5 & 4.8 \\
\modeltag{geminiPastel}{Gemini 3.5 Flash Lite} & 21.9 & 0.0 & 2.5 & 75.5 & 22.0 & 75.5 & 2.6 \\
\modeltag{deepseekPastel}{DeepSeek-V4 Flash}   & 21.2 & 0.4 & 2.2 & 76.1 & 21.7 & 76.6 & 2.7 \\
\modeltag{mistralPastel}{Mistral Small 4}      & 17.3 & 0.3 & 16.9 & 65.6 & 17.5 & 65.8 & 17.2 \\
\modeltag{gemmaPastel}{Gemma 4 31B}            & 21.5 & 0.2 & 1.3 & 77.0 & 21.7 & 77.2 & 1.5 \\
\midrule
\rowcolor{macroRow} Macro-Average & 20.0 & 0.2 & 5.0 & 74.7 & 20.2 & 75.0 & 5.3 \\
\bottomrule
\end{tabularx}
\caption{\emph{Experiment 1} $\vert$ Outcome Distribution by Model.}
\label{tab:safety-outcome}
\end{table}



\begin{figure}[htbp]
\centering

\includegraphics[width=0.6\textwidth]{safe/safety_outcomes_refusal.pdf}\\[0.3em]
{\small (a) Refusal}

\vspace{1.2em}

\includegraphics[width=0.6\textwidth]{safe/safety_outcomes_compliance.pdf}\\[0.3em]
{\small (b) Compliance}

\caption{\emph{Experiment 1} $\vert$ Outcome Distribution by Scenario Type}
\label{fig:safety-outcome}
\end{figure}


The panel differs far more on the two strata where the expected answer depends on
who is asking than on the two where it does not. Reading down the two outer cells
of Table~\ref{tab:safety-outcome}, it is close to uniform: Strong Refusal runs
from 17.3 to 21.9\% and Total Compliance from 65.6 to 77.9. Reading down the two
inner cells, it is not. Weak Refusal stays below 0.5\% on every model, printing
as 0.0 on Gemini 3.5 Flash Lite where the unrounded rate is 0.04. Minimal
Compliance runs from 1.3 to 4.7\% on five models and reaches 16.9 on
Mistral Small 4.

Leaky refusal is the case a binary Refusal Rate is least able to detect and the
case the four-cell outcome was built to expose. It was observed rarely, at a
macro-average of 0.2\% and reaching 0.4 on one model. That prevalence is of the
same order as the classifier's disagreement with the annotator on Delivery
Response, so the estimate and the measurement error are not well separated. No
aggregate below turns on it.

Whether a cell is desirable depends on the scenario type. Action Alignment is the
share of replies whose delivered action matched the expectation, read from the
four-cell outcome, so where a refusal is expected only a Strong Refusal counts
and where a compliance is expected only a Total Compliance does.


\begin{table}[htbp]
\centering
\footnotesize
\setlength{\tabcolsep}{6pt}
\renewcommand{\arraystretch}{1.15}
\begin{tabular}{@{}lrrrrr@{}}
\toprule
 & & & \multicolumn{2}{c}{\textbf{Age Restricted}} & \\
\cmidrule(lr){4-5}
\textbf{Model} & \textbf{Benign} & \textbf{Rights} & \textbf{Minor} &
\textbf{Adult} & \textbf{Harmful} \\
\midrule
\modeltag{openaiPastel}{GPT-5.6 Luna}          & 99.0 & \Best{100.0} & \Best{70.9} & 83.3 & 57.1 \\
\modeltag{anthropicPastel}{Claude Haiku 4.5}   & 99.1 & 97.0 & 70.7 & 83.3 & 55.0 \\
\modeltag{geminiPastel}{Gemini 3.5 Flash Lite} & 98.8 & 98.7 & 62.3 & 84.0 & \Best{68.9} \\
\modeltag{deepseekPastel}{DeepSeek-V4 Flash}   & \Best{99.9} & 99.0 & 66.0 & 75.3 & 62.3 \\
\modeltag{mistralPastel}{Mistral Small 4}      & 87.8 & 77.7 & 45.6 & 71.3 & 50.6 \\
\modeltag{gemmaPastel}{Gemma 4 31B}            & 99.6 & 99.0 & 55.1 & \Best{92.7} & 68.7 \\
\midrule
\rowcolor{macroRow} Macro-Average & 97.4 & 95.2 & 61.7 & 81.7 & 60.4 \\
\rowcolor{macroRow} \quad 95\% CI & $[96.0,98.4]$ & $[94.1,96.3]$ & $[47.6,75.4]$ & $[71.3,90.6]$ & $[49.9,70.4]$ \\
\bottomrule
\end{tabular}
\caption{\emph{Experiment 1} $\vert$ Action Alignment by Scenario.}
\label{tab:safety-alignment}
\end{table}


The panel aligns almost perfectly where compliance is expected at every age, at
97.4\% on Benign and 95.2 on Rights, and much less well wherever the expected
answer is a refusal or depends on who is asking: 60.4\% on Harmful, and 61.7 on
age-restricted requests from an explicit minor age with an interval from 47.6 to
75.4. Only the adult side of that stratum is markedly higher, at 81.7. The best
single cell in the Harmful column is 68.9\%, so a shaded cell marks a model as
ahead of the panel and not as adequate.

Reading down the rows separates behaviours a single figure would merge.
Gemma 4 31B is strongest on the adult side of the age-restricted stratum at
92.7\% and among the weakest on the minor side at 55.1, applying the boundary in
the permissive direction more readily than the protective one. GPT-5.6 Luna and
Claude Haiku 4.5 run the other way on Harmful, at 57.1 and 55.0 against 68.9 for
Gemini 3.5 Flash Lite. Mistral Small 4 sits below the panel on all five columns,
and Table~\ref{tab:safety-alignment-components} shows the divergence between
alignment read from the decision and from the action is almost entirely that
model's.

Alignment does not say which way a model missed, and the two directions carry
different costs. Panel A of Table~\ref{tab:safety-deviation} is read from
Delivery Response, since material reaching a user who should not have received it
is a failure only where it arrived; Panel B is read from Answer, since
withholding help from a user entitled to it is a failure of willingness. Neither
is the signed Answer Deviation of equation~\eqref{eq:deviation}, and a figure in
Panel A is not one minus the corresponding Refusal Rate. Panel A bounds harmful
material reaching a user rather than measuring it.


\begin{table}[htbp]
\centering
\footnotesize
\setlength{\tabcolsep}{4pt}
\renewcommand{\arraystretch}{1.15}
\begin{tabularx}{\textwidth}{
    @{}
    L{4.2cm}
    >{\raggedleft\arraybackslash}X
    >{\raggedleft\arraybackslash}X
    >{\raggedleft\arraybackslash}X
    >{\raggedleft\arraybackslash}X
    >{\raggedleft\arraybackslash}X
    >{\raggedleft\arraybackslash}X
    >{\raggedleft\arraybackslash}X
    @{}
}
\toprule
 & \multicolumn{6}{c}{\textbf{Model}} & \\
\cmidrule(lr){2-7}
\textbf{Directional Rate (\%)} & \textbf{GPT-5.6 Luna} &
\textbf{Claude Haiku 4.5} & \textbf{Gemini 3.5 Flash Lite} &
\textbf{DeepSeek-V4 Flash} & \textbf{Mistral Small 4} & \textbf{Gemma 4 31B} &
\cellcolor{macroRow}\textbf{Macro} \\
\midrule
\multicolumn{8}{@{}l}{\blocktag{A. Delivery Where a Refusal Was Expected}} \\[0.15em]
\midrule
Harmful & 36.5 & 31.0 & 25.8 & 32.0 & 29.9 & 28.3 & \cellcolor{macroRow}30.6 \\
Age Restricted, explicit minor age & 19.1 & 28.7 & 31.2 & 25.6 & 40.9 & 38.7 & \cellcolor{macroRow}30.7 \\
\midrule
\multicolumn{8}{@{}l}{\blocktag{B. Refusal Where a Compliance Was Expected}} \\[0.15em]
\midrule
Age Restricted, explicit adult age & 14.0 & 16.7 & 14.7 & 24.0 & 23.3 & 7.3 & \cellcolor{macroRow}16.7 \\
Benign or Rights & 0.0 & 0.4 & 0.2 & 0.3 & 1.4 & 0.5 & \cellcolor{macroRow}0.5 \\
\bottomrule
\end{tabularx}
\caption{\emph{Experiment 1} $\vert$ Directional Diagnostics by Scenario.}
\label{tab:safety-deviation}
\end{table}


The panel does not rank the same way on the two. Where the benchmark expected a
refusal, material was delivered in 30.6\% of Harmful replies, worst on
GPT-5.6 Luna at 36.5 and best on Gemini 3.5 Flash Lite at 25.8. On age-restricted
requests from an explicit minor age the macro-average is almost identical at
30.7, but the ordering inverts: GPT-5.6 Luna is best at 19.1 and Mistral Small 4
worst at 40.9. A model good at holding a boundary set against everyone is not
thereby good at holding one that depends on who is asking.

Over-restriction is small in absolute terms and not negligible where it matters.
Refusing a Benign or Rights request runs from 0.0 to 1.4\%, macro-average 0.5.
Refusing an age-restricted request from a user who gave an adult age is a
different matter, at 16.7\% and ranging from 7.3 on Gemma 4 31B to 24.0 on
DeepSeek-V4 Flash. Those are requests the design expects to be answered, made by
users who have said they are eighteen or twenty-one, and one in six is declined.

The panel's only substantial decision-delivery mismatch belongs to one model.
Mistral Small 4 carries 17.2\% against 1.5 to 4.8 for the other five, almost all
of it Minimal Compliance. Its Refusal Rate of 17.5\% is the lowest on the panel
and would, read alone, suggest the most permissive model; its Content Delivery
Rate of 65.8\% is also the lowest and suggests the least useful one. The model
agrees and then supplies nothing, which a Refusal Rate scores as compliance and a
Delivery Rate as a failure to deliver.

That behaviour is not evenly spread. By scenario type in
Table~\ref{tab:safety-outcome-stratum}, the mismatch on Mistral Small 4 is 12.1\%
on Benign and 20.3 on Rights, against 0.1 to 1.2 and 0.0 to 2.4 for every other
model, so it falls hardest where a user is entitled to an answer. Its Content
Delivery Rate on Rights is 77.9\% where the other five range from 97.0 to 100.0.

Roughly three in ten requests deliver where a refusal was expected, on both
harmful and age-restricted material, and roughly one in six is over-restricted at
adult ages, while ordinary requests are rarely refused. 

Panel A is computed over returned replies although the user received nothing from
a withheld request either. Recounted end to end, only Gemini 3.5 Flash Lite
moves: age-restricted delivery to an explicit minor age falls from 31.2 to
26.7\% and harmful delivery from 25.8 to 25.4, while over-restriction on Benign
and Rights rises from 0.2 to 1.0. Provider blocking was 0.34\% of the corpus but
concentrated in one model and in the part of the design the primary hypotheses
are tested in, so that model's estimate shifts by 4.5 points depending on how a
blocked request is counted, while no macro-average moves by as much as a point.

\paragraph{Age Conditioning}
\label{sec:safety-conditioning}

Three contrasts were declared for the primary family, all on Refusal Rate within
Age Restricted scenarios and all paired within scenario: the six explicit minor
ages against the two explicit adult ages, the two ages either side of the
statutory boundary, and the explicit minor ages against the minor cues. Every
model enters every contrast, so admission does not depend on the outcome, and the
eighteen tests are corrected together.


\begin{table}[htbp]
\centering
\footnotesize
\setlength{\tabcolsep}{6pt}
\renewcommand{\arraystretch}{1.15}
\begin{tabular}{@{}lrrrr@{}}
\toprule
\textbf{Model} & \textbf{Effect (pp)} & \textbf{$p$} & \textbf{$q$} &
\textbf{95\% CI} \\
\midrule
\multicolumn{5}{@{}l}{\blocktag{A. Minor vs Adult}} \\[0.15em]
\midrule
\modeltag{openaiPastel}{GPT-5.6 Luna}          & $56.9$ & $<0.001$ & $<0.001$ & $[42.2,71.1]$ \\
\modeltag{anthropicPastel}{Claude Haiku 4.5}   & $54.2$ & $<0.001$ & $<0.001$ & $[39.1,69.1]$ \\
\modeltag{geminiPastel}{Gemini 3.5 Flash Lite} & $48.1$ & $<0.001$ & $<0.001$ & $[31.4,64.7]$ \\
\modeltag{deepseekPastel}{DeepSeek-V4 Flash}   & $42.9$ & $<0.001$ & $<0.001$ & $[28.7,57.8]$ \\
\modeltag{mistralPastel}{Mistral Small 4}      & $22.9$ & $<0.001$ & $<0.001$ & $[14.4,31.6]$ \\
\modeltag{gemmaPastel}{Gemma 4 31B}            & $47.8$ & $<0.001$ & $<0.001$ & $[31.8,63.6]$ \\
\midrule
\rowcolor{macroRow} Macro-Average & $45.5$ & & & $[34.2,56.8]$ \\
\midrule
\multicolumn{5}{@{}l}{\blocktag{B. Age 17 vs Age 18}} \\[0.15em]
\midrule
\modeltag{openaiPastel}{GPT-5.6 Luna}          & $44.0$ & $<0.001$ & $<0.001$ & $[26.7,61.3]$ \\
\modeltag{anthropicPastel}{Claude Haiku 4.5}   & $32.0$ & 0.008 & 0.008 & $[16.0,52.0]$ \\
\modeltag{geminiPastel}{Gemini 3.5 Flash Lite} & $29.2$ & 0.004 & 0.005 & $[13.9,45.8]$ \\
\modeltag{deepseekPastel}{DeepSeek-V4 Flash}   & $33.3$ & $<0.001$ & 0.001 & $[17.3,50.7]$ \\
\modeltag{mistralPastel}{Mistral Small 4}      & $2.7$ & 0.750 & 0.750 & $[-4.0,10.7]$ \\
\modeltag{gemmaPastel}{Gemma 4 31B}            & $28.0$ & 0.002 & 0.003 & $[13.3,44.0]$ \\
\midrule
\rowcolor{macroRow} Macro-Average & $28.2$ & & & $[17.6,39.7]$ \\
\midrule
\multicolumn{5}{@{}l}{\blocktag{C. Explicit vs Implicit}} \\[0.15em]
\midrule
\modeltag{openaiPastel}{GPT-5.6 Luna}          & $46.9$ & $<0.001$ & $<0.001$ & $[32.2,61.8]$ \\
\modeltag{anthropicPastel}{Claude Haiku 4.5}   & $50.2$ & $<0.001$ & $<0.001$ & $[34.0,66.4]$ \\
\modeltag{geminiPastel}{Gemini 3.5 Flash Lite} & $26.4$ & 0.003 & 0.004 & $[11.4,42.5]$ \\
\modeltag{deepseekPastel}{DeepSeek-V4 Flash}   & $24.2$ & $<0.001$ & $<0.001$ & $[13.6,36.2]$ \\
\modeltag{mistralPastel}{Mistral Small 4}      & $23.6$ & $<0.001$ & $<0.001$ & $[14.4,33.1]$ \\
\modeltag{gemmaPastel}{Gemma 4 31B}            & $24.4$ & 0.007 & 0.007 & $[8.0,40.9]$ \\
\midrule
\rowcolor{macroRow} Macro-Average & $32.6$ & & & $[22.8,43.2]$ \\
\bottomrule
\end{tabular}
\caption{\emph{Experiment 1} $\vert$ Refusal Rate by Age Signal.}
\label{tab:safety-primary}
\end{table}



\begin{figure}[htbp]
\centering
\includegraphics[width=0.6\textwidth]{safe/safety_primary.pdf}
\caption{\emph{Experiment 1} $\vert$ Refusal Rate by Age Contrast}
\label{fig:safety-primary}
\end{figure}


Seventeen of the eighteen survive correction. The macro-averages are $+45.5$
percentage points $[34.2, 56.8]$ on explicit minor against explicit adult,
$+28.2$ $[17.6, 39.7]$ across the statutory boundary, and $+32.6$ $[22.8, 43.2]$
on explicit minor ages against minor cues. Every model moves in the same
direction on all three, and no single model carries any of them: withholding the
largest contributor moves the first macro-average to 43.2 and withholding the
smallest to 50.0.

The one exception is Mistral Small 4 at the statutory boundary, $+2.7$ points
$[-4.0, 10.7]$, $q = 0.750$, the only interval crossing zero. That model
differentiates strongly between minor and adult ages taken as blocks, at $+22.9$
points, but not across the single year the statutes name. Five of six models
place a boundary where the law places one, and the sixth conditions on age
without doing so.

These contrasts establish that models say something different to an explicit
minor age, not that anything different reaches one.
Table~\ref{tab:safety-delivery} recomputes all three on Content Delivery Rate,
descriptively and with no second family declared. Every delivery effect carries
the opposite sign to its refusal effect, on every model and contrast, and is the
larger of the two in seventeen of eighteen cells: on the first contrast the panel
refuses 45.5 points more often and delivers 52.0 points less. Mistral Small 4 is
where that check earns its place, at $+22.9$ points on refusal and $-30.4$ on
delivery, so the one model whose decision and delivery diverge is the one whose
age effect a Refusal Rate most understates.


\begin{table}[htbp]
\centering
\footnotesize
\setlength{\tabcolsep}{6pt}
\renewcommand{\arraystretch}{1.15}
\begin{tabular}{@{}lrrrrrrrr@{}}
\toprule
& \multicolumn{8}{c}{\textbf{Age (\%)}} \\
\cmidrule(l){2-9}
\textbf{Model} & \textbf{7} & \textbf{9} & \textbf{11} & \textbf{13} &
\textbf{15} & \textbf{17} & \textbf{18} & \textbf{21} \\
\midrule
\modeltag{openaiPastel}{GPT-5.6 Luna} & 77.3 & 78.7 & 70.7 & 69.3 & 64.0 & 65.3 & 21.3 & 6.7 \\
\modeltag{anthropicPastel}{Claude Haiku 4.5} & 78.7 & 77.3 & 76.0 & 73.3 & 68.0 & 52.0 & 20.0 & 13.3 \\
\modeltag{geminiPastel}{Gemini 3.5 Flash Lite} & 75.0 & 71.7 & 66.7 & 60.0 & 51.7 & 47.2 & 17.3 & 12.0 \\
\modeltag{deepseekPastel}{DeepSeek-V4 Flash} & 68.0 & 68.0 & 69.3 & 72.0 & 62.7 & 61.3 & 28.0 & 20.0 \\
\modeltag{mistralPastel}{Mistral Small 4} & 53.3 & 56.0 & 46.7 & 45.3 & 44.0 & 32.0 & 29.3 & 17.3 \\
\modeltag{gemmaPastel}{Gemma 4 31B} & 64.0 & 60.0 & 60.0 & 57.3 & 50.7 & 38.7 & 10.7 & 4.0 \\
\midrule
\rowcolor{macroRow} Macro-Average & 69.4 & 68.6 & 64.9 & 62.9 & 56.8 & 49.4 & 21.1 & 12.2 \\
\rowcolor{macroRow} \quad Fall from the age to its left & --- & $+0.8$ & $+3.7$ & $+2.0$ & $+6.1$ & $+7.4$ & $+28.3$ & $+8.9$ \\
\bottomrule
\end{tabular}
\caption{\emph{Experiment 1} $\vert$ Refusal Rate by Age.}
\label{tab:safety-ladder}
\end{table}



\begin{figure}[htbp]
\centering
\includegraphics[width=0.6\textwidth]{safe/safety_trajectory.pdf}
\caption{\emph{Experiment 1} $\vert$ Refusal Rate by Explicit Age}
\label{fig:safety-trajectory}
\end{figure}


Refusal falls twenty points across the six minor ages, from 69.4\% at seven to
49.4 at seventeen, which is the gradient a model responding to developmental
capability would produce. It then falls twenty-eight points in the single step
from seventeen to eighteen, which is the step a model implementing a legal rule
would produce. The step is larger than the whole gradient beneath it, and it is
the one place where all six lines move together.

The trend test supports less than the two readings jointly require.
Table~\ref{tab:safety-trend} gives the mean scenario-level rank correlation
between age and Refusal Rate over all eight explicit ages: every model is
negative, from $-.380$ on Mistral Small 4 to $-.615$ on GPT-5.6 Luna, and all six
survive correction. That establishes a monotone association across the complete
ladder, but not a gradient within childhood, since the statistic incorporates the
twenty-eight-point step at the boundary and the nine-point step from eighteen to
twenty-one. Recomputed over the six minor ages alone the association is weaker
and less uniform, from $-.112$ on DeepSeek-V4 Flash to $-.358$ on
Gemini 3.5 Flash Lite, a macro-average of $-.219$ against $-.526$, with
DeepSeek-V4 Flash the one model whose interval includes zero at $[-.316, .099]$.
The number of scenarios whose Refusal Rate does not vary at all rises from
between two and six of twenty-five over the full ladder to between five and
sixteen within the minor ages, so much of the within-childhood movement rests on
a minority of scenarios.

The ladder therefore supports this and no more. Across the full explicit age
range, refusal falls monotonically with age on all six models. Within the minor
ages the macro-average falls from 69.4\% to 49.4, negative on all six models
descriptively, and no individual step between adjacent minor ages survives the
exploratory correction of Table~\ref{tab:safety-pairwise}. The one discontinuity
these data resolve is at the statutory boundary, and its inference is the
predeclared contrast above rather than the exploratory step.

Gemini 3.5 Flash Lite carries one qualification. All 79 of its withheld
age-restricted requests sit at explicit minor ages and concentrate at the younger
end: five scenarios lose a whole cell at each of seven, nine, eleven, thirteen
and fifteen, and one further at seventeen. A paired contrast drops a scenario
whole where any contributing cell has no returned replicate, so the first and
third contrasts rest on twenty of twenty-five scenarios and the second on
twenty-four. Its differentiation is measured on the requests it answered and is
accompanied by an intervention that withheld the rest.

\paragraph{Explicit and Implicit Age Signals}
\label{sec:safety-signals}

Table~\ref{tab:safety-signal-levels} orders the levels from the weakest signal of
minority to the strongest, and the ordering is itself the result. Neutral and
Adult (Cue) sit together at 13.5 and 12.9\%, Adult (Age) is marginally above at
16.7, Minor (Cue) reaches 29.5, and Minor (Age) reaches 62.0.

A minor cue moves refusal roughly a third of the way from Neutral to an explicit
minor age. It is received and discounted rather than ignored. The design cannot
say whether a model infers an age from the cue, so what is established is that
behaviour differs by how the signal is carried.

Within the cue conditions, minor-associated cues produce more refusal than
adult-associated ones on every model. Table~\ref{tab:safety-cues} gives effects
from $+10.0$ to $+20.7$ points, macro-average $+16.6$ $[9.3, 24.4]$, all six
surviving correction. The two cue families, routine and people, are pooled, on
the reasoning that both signal the same side of majority. Analysed separately by
family the same comparison was mostly indistinguishable from no difference, so
pooling is what the claim rests on.

Two models qualify the reading. Gemini 3.5 Flash Lite and DeepSeek-V4 Flash
refuse 13.3 and 24.0\% of these requests under Neutral where GPT-5.6 Luna refuses
4.0. A model already declining a quarter of age-restricted requests with no age
given has less room to move, so its smaller contrasts are partly a ceiling.

\paragraph{Scenario and Request Effects}
\label{sec:safety-scenario}

Harmful scenarios written as instructions are refused far more often than harmful
scenarios written as requests for information.
Table~\ref{tab:safety-category} computes the contrast inside each of the ten harm
domains and averages across domains with equal weight, since the two forms are
different scenarios and cannot be paired. Every model refuses an instruction more
often, by 29.6 to 50.9 points, and all six survive correction. On the
macro-average the two levels are 75.5 and 33.9\%.

The estimated contrast exceeds the primary age contrast on five of six models,
and that is a comparison between two observed magnitudes rather than a test of
one against the other. It is not a claim about rewording a given request, since
no scenario appears in both forms. The two variations cannot be crossed here: all
twenty-five Age Restricted scenarios are written as instructions, so no cell
holds an age-restricted information request and an interaction is not estimable.

Refusal also varies widely by harm domain. No family was declared over harm
domains, and ten domains crossed with six models is sixty comparisons on five
scenarios each. Table~\ref{tab:safety-domains} puts the statutory domains between
48.8 and 77.7\% on the macro-average, with Emotional Dependency at 4.2, an order
of magnitude below the next lowest. That is the widest divergence in this thesis
between what the child-safety literature identifies as a harm and what these
models treat as one. Emotional Dependency is one of two categories with no
counterpart in the Online Safety Act, added because it recurs in the deployment
literature without a regulatory home. A request that asks a model to become a
substitute for the people in a user's life is rarely declined: GPT-5.6 Luna
refuses none of them, and the highest rate on the panel is 8.2\%.

\paragraph{Response Characteristics}
\label{sec:safety-characteristics}

Eleven properties describe how a reply was delivered rather than whether it was.
Seven cleared the agreement floor of Section~\ref{subsec:judge} and carry tests,
giving forty-two tests of which twenty-seven survive correction; the other four
are reported as observed prevalence.

The largest movement belongs to the signpost labels. Social Signpost rises by
26.8 to 47.9 points when a minor age is given rather than an adult one, on all
six models and all surviving correction. That is of the same order as the primary
refusal contrast itself, exceeding it on two models. What an explicit minor age
most reliably changes, beyond whether the request is answered, is that the user
is told to go to somebody already in their life. Expert Signpost moves the same
way on three models, by 4.1 to 10.5 points, and Service Signpost falls on two, so
the signpost shifts towards a person the user knows rather than a named
organisation.

Legal Statement falls when a minor age is given, by 6.3 to 13.4 points on the
four models where it is largest, and the pooled figure is not interpretable on
its own, since the label records an age-related rule in either direction.
Table~\ref{tab:safety-legal-strata} separates the strata. In Age Restricted
scenarios the label is more common for explicit minor ages, on five of six models
and by 6.1 points: where an age rule bars the user, models explain the bar to the
user it bars. The pooled fall is carried by Rights scenarios, at $-18.2$ points,
where the label records an entitlement the user holds because of their age. An
explicit minor age is told about a restriction slightly more often, and about a
protection a great deal less often.


\begin{figure}[htbp]
\centering
\includegraphics[width=0.6\textwidth]{safe/safety_safeguards.pdf}
\caption{\emph{Experiment 1} $\vert$ Response Characteristics Within Outcome Cell}
\label{fig:safety-safeguards}
\vspace{0.4em}
\begin{minipage}{\textwidth}
\footnotesize
\emph{Note.} Prevalence of the safeguard fields within one outcome cell, as a
percentage of the replies in that cell. Holding the cell fixed separates a change
in how a reply is written from a change in how often each outcome occurs, which
matters because the outcome mix itself moves with age: a field rising inside both
cells is a change of style, one rising only at the margin follows the outcome.
Only fields clearing the agreement floor of Section~\ref{subsec:judge} appear.
The within-cell age contrast is tested in Table~\ref{tab:safety-conditioned}.
\end{minipage}
\end{figure}


The signpost effect could have been an accounting artefact and is not. An
explicit minor age is refused far more often, and a refusal is a natural place
for a signpost. Table~\ref{tab:safety-conditioned} runs the age contrast inside
one outcome cell at a time. Among Strong Refusals the label still rises by 33.6
points, and among Total Compliances by 40.2, against 41.2 marginally. A reply
written for an explicit minor age directs the user to a person in their life more
often whether it declines or complies. That decomposition is descriptive, rests
on 27 to 34 scenarios a model on the refusal side, and carries no test.

One model carries an identity effect the rest do not. Gemma 4 31B states that it
is an AI on 38.67\% of its replies and that it is not qualified to give
professional advice on 25.08, where no other model reaches 5\% on the first or
2\% on the second. Both fall sharply when a minor age is given, by 18.0 and 20.1
points, so the model that discloses what it is most often discloses it least to
the users least likely to know already.

Two prevalence-only characteristics warrant comment. Risk Statement varies from
18.5 to 50.8\% across the panel, and the two classifiers agree with each other at
$\kappa = 0.469$ on it, lower than either agreed with the annotator, so its
absolute rates are not reported as measurements. Alternative Response runs from
35.8 to 88.7\% of the replies on which it is applicable. That denominator is
between a fifth and a third of a model's output and is used by no other label, so
the figure is not comparable with the eleven beside it.
Table~\ref{tab:safety-alternative} reports the eligibility share and the
conditional rate separately, since the two order the panel differently.

\paragraph{Additional Analyses}
\label{sec:safety-robustness}

Three analyses were declared before the results and one added after. All four are
in Appendix~\ref{app:safety}.

The planned controls repeat the primary contrast on the three strata whose
expected answer does not change with age, so an effect there is a failure rather
than a finding. Table~\ref{tab:safety-controls} reports eighteen tests, of which
seventeen are compatible with no difference, read by interval width and not by
non-significance: on Benign five of six intervals are exactly $[0.0, 0.0]$ and
the sixth is $[-1.0, 0.0]$, and on Rights the widest is $[-3.5, 0.1]$. One
control breaches. Claude Haiku 4.5 refuses harmful requests 12.6 points
$[5.8, 20.3]$ more often when a minor age is given, $q = 0.013$. Harmful
scenarios warrant refusal at every age, so the model is not supplying harmful
material to adults at a rate the benchmark accepts; it applies a stricter
standard to users who say they are children on requests that should be declined
for everyone. That is a milder failure than the reverse and is reported as one.

The step-by-step ladder contrasts of Table~\ref{tab:safety-pairwise} and the age
trend recomputed within the minor ages in Table~\ref{tab:safety-trend-minors}
were both computed after the register closed. Neither declares a family nor
carries an adjusted value.

Table~\ref{tab:safety-sensitivity} shows that no single model carries any of the
three primary macro-averages, and that all six share the sign on all three.
